% Study sheet: Volume I, Chapter 11, Hardware Acceleration.
% Numbers: \val{...} from hw_acceleration_studysheet_numbers.tex, written by hw_acceleration_studysheet_numbers.py.
\documentclass{studysheet}
\input{hw_acceleration_studysheet_numbers}
\sheetvolume{1}
\sheetchapter{11}{Hardware Acceleration}
\sheetsources{Numbers from the mlsysim registry and the chapter's scenarios (H100 dense FP16 and INT8 paths; 45 nm energies); assumed inputs are marked.}

\begin{document}
% ===================================================================== FRONT
\frontside
\begin{bigidea}
Arithmetic is cheap; \hl{moving bytes is not}. An accelerator wins by keeping operands close and
reusing them, so it reaches its peak only on work that does enough math per byte. Otherwise memory
bandwidth, or work the accelerator never touches, sets the time.
\end{bigidea}

\begin{leftcol}
\sheetsec{Key ideas}
\begin{ideas}
\idea{Moving data costs more than math.}{A DRAM access costs about \val{dramovermul}$\times$ the energy of
  an FP32 multiply and hundreds of cycles. Specialized hardware exists to avoid that trip.}
\idea{Reuse is the design.}{Matrix units reuse each fetched operand across many multiply-accumulates;
  a systolic array passes it cell to cell. Vector units take element-wise work and reductions.}
\idea{A peak belongs to one path.}{INT8 doubles the H100's compute ceiling, not its bandwidth, so the
  ridge doubles. Zeros are skipped only in the supported 2:4 pattern.}
\idea{Batch one streams every weight.}{Decoding a token reads all weights for 2 operations each. A dense
  layer reuses each weight $\BatchSize$ times, so $\ArithmeticIntensity \approx \BatchSize$ at FP16.}
\idea{Mapping decides the bytes.}{Keeping weights or partial sums local (weight- or output-stationary),
  tiles that fit on chip, and fused kernels cut HBM traffic: the kernel moves right on the roofline.}
\idea{Profile before buying.}{The roofline bounds a kernel; Amdahl bounds the job. A profile shows
  which limit you actually hit.}
\end{ideas}


\vspace{12pt}
\sheetsec{What limits it, and which lever}
\begin{sheettable}{@{}L L@{}}
\thead{You see} & \thead{Lever}\\ \midrule
Below the ridge; HBM near peak, kernels back to back & dense layer ($\ArithmeticIntensity \approx \BatchSize$): bigger batch, narrower weights\\
Element-wise or normalization, $\ArithmeticIntensity$ near 1--2 & fuse into a neighbor: batching does not help\\
$\ArithmeticIntensity$ above the ridge & shapes and tiles that keep the matrix units full\\
Low HBM traffic, gaps between short kernels & launch overhead: compile, fuse, capture the graph\\
Job speedup far below the kernel gain & Amdahl: fix the unaccelerated work\\
\end{sheettable}
\end{leftcol}%
\begin{rightcol}
\sheetsec{Formulas you will use}
\fcard{Roofline and ridge}%
  {\AttainableComputeRate = \min(\PeakRate,\ \Bandwidth \times \ArithmeticIntensity), \qquad \RooflineRidgePoint = \frac{\PeakRate}{\Bandwidth}}%
  {Intensity $\ArithmeticIntensity$ = operations per byte moved. Below the ridge $\RooflineRidgePoint$, bandwidth
   $\Bandwidth$ sets the rate; above it, the peak $\PeakRate$ of that precision path does.}%
  {Bounds at HBM. Batch-one decode: time per token $\ge$ weight bytes $\div \Bandwidth$. H100: \val{hpeaksixteen} TFLOP/s FP16,
   \val{hpeakeight} TOPS INT8, \val{hbw} TB/s.}%
  {\begin{tikzpicture}[sheet]
     % x: log10(I) -0.3..3.7 -> 30..202 ; y: log10(TFLOP/s) 0..3.5 -> 0..105
     \pgfmathsetmacro\sx{43}\pgfmathsetmacro\sy{30}
     \foreach \tk/\tl in {1/1,10/10,100/100,1000/1000}{
       \pgfmathsetmacro\x{30+(log10(\tk)+0.3)*\sx}
       \draw[gridline] (\x,0) -- (\x,105);
       \node[small, anchor=north] at (\x,-3) {\tl};}
     \foreach \tk/\tl in {1/1,10/10,100/100,1000/1000}{
       \pgfmathsetmacro\y{log10(\tk)*\sy}
       \draw[gridline] (30,\y) -- (202,\y);
       \node[small, anchor=east] at (27,\y) {\tl};}
     \node[small, anchor=north east] at (202,-12) {operations per byte};
     \node[small, anchor=south, rotate=90] at (8,52) {tera-ops/s};
     \pgfmathsetmacro\bw{\val{hbw}}
     \pgfmathsetmacro\rs{\val{ridgesixteen}}\pgfmathsetmacro\re{\val{ridgeeight}}
     \pgfmathsetmacro\xs{30+(log10(\rs)+0.3)*\sx}\pgfmathsetmacro\xe{30+(log10(\re)+0.3)*\sx}
     \pgfmathsetmacro\ys{log10(\val{hpeaksixteen})*\sy}
     \pgfmathsetmacro\yeight{log10(\re*\bw)*\sy}
     \pgfmathsetmacro\ylow{(log10(\bw)-0.3)*\sy}
     \draw[accent, line width=1.1pt] (30,\ylow) -- (\xs,\ys) -- (202,\ys);
     \draw[accent, dashed, line width=0.7pt] (\xs,\ys) -- (\xe,\yeight) -- (202,\yeight);
     \draw[dim, dashed] (\xs,0) -- (\xs,\ys);
     \node[note, anchor=south east] at (\xs-2,3) {ridge \val{ridgesixteen}};
     \node[note, anchor=south east] at (\xe-1,\yeight+2) {ridge \val{ridgeeight}};
     \node[lab, anchor=west] at (205,\ys) {FP16};
     \node[lab, anchor=west] at (205,\yeight+2) {INT8};
     \pgfmathsetmacro\xg{30+(log10(\val{gintensity})+0.3)*\sx}\pgfmathsetmacro\yg{log10(\val{gattain})*\sy}
     \filldraw[fill=accent, draw=accent] (\xg,\yg) circle (2.4pt);
     \node[note, anchor=west] at (\xg+5,\yg-4) {GPT-2 XL decode};
     \pgfmathsetmacro\xc{30+(log10(\val{convsixteen})+0.3)*\sx}
     \filldraw[fill=accent, draw=accent] (\xc,\ys) circle (2.4pt);
     \node[note, anchor=north west] at (\xe+4,\ys-5) {ResNet-50 conv};
   \end{tikzpicture}}
\fcard{Amdahl's law}%
  {\text{Speedup} = \frac{1}{(1-\ParallelFraction) + \ParallelFraction/\AcceleratorGain} \;\le\; \frac{1}{1-\ParallelFraction}}%
  {$\ParallelFraction$: share of baseline time the faster path covers; $\AcceleratorGain$: that path's gain.}%
  {Data loading, host work, transfers and launches stay in $1-\ParallelFraction$.}{}
\vspace{10pt}
\sheetsec{Energy per 32-bit event (45 nm)}
{\centering
\begin{tikzpicture}[sheet]
  \draw[dim] (10,0) -- (240,0);
  \foreach \tk/\tl in {0.1/0.1,1/1,10/10,100/100,1000/1000 pJ}{
    \pgfmathsetmacro\x{10+(log10(\tk)+1)*57.5}
    \draw[dim] (\x,0) -- (\x,-2.5);
    \node[small, anchor=north] at (\x,-4) {\tl};}
  \pgfmathsetmacro\xm{10+(log10(\val{emulfp})+1)*57.5}
  \pgfmathsetmacro\xs{10+(log10(\val{esram})+1)*57.5}
  \pgfmathsetmacro\xd{10+(log10(\val{edram})+1)*57.5}
  \foreach \xx in {\xm,\xs}{\filldraw[fill=accentsoft, draw=accent, line width=0.5pt] (\xx,0) circle (2.4pt);}
  \filldraw[fill=accent, draw=accent] (\xd,0) circle (2.4pt);
  \node[note, anchor=south east] at (\xm+1,5) {FP32 multiply \val{emulfp}};
  \node[note, anchor=south west] at (\xs-1,15) {8 KB SRAM read \val{esram}};
  \draw[dim] (\xs,3) -- (\xs,14);
  \node[note, anchor=south] at (\xd-6,5) {DRAM \val{edram}};
\end{tikzpicture}\par}
\end{rightcol}

% ===================================================================== BACK
\backside
\begin{answers}
\ans{Finish}{\val{gops} G $\div$ \val{yweightgb} GB $=$ \val{yintensity} op/byte, far below \val{ridgeeight}: bandwidth
  bound. Floor \val{yweightgb} GB $\div$ \val{hbw} TB/s $=$ \val{yfloorms} ms, a floor \val{yspeedup}$\times$ lower
  (\val{ycut}\% less time), all from halved bytes; the doubled INT8 peak is never reached.}
\ans{1}{One over (\val{qoneserialfrac} $+$ \val{qonepg}) $=$ \val{qonespeed}$\times$; ceiling one over \val{qoneserialfrac}
  $=$ \val{qoneceiling}$\times$. Fix the data loading first: it caps the job.}
\ans{2}{\val{edram} $\div$ \val{esram} $=$ \val{dramoversram} SRAM reads. Streaming: \val{qtwonaiveacc} $\times$ \val{edram} $+$
  \val{qtwomac} $=$ \val{qtwonaive} pJ per MAC; array: \val{qtwosysacc} fetches per \val{qtwowidth} MACs add \val{qtwodram} pJ,
  \val{qtwosys} pJ in all, \val{qtworatio}$\times$ less. Reuse, not faster math.}
\columnbreak
\ans{3}{All against ridge \val{ridgesixteen}. LayerNorm \val{qthreeln}: $\le$ \val{qthreelnrate} TFLOP/s, fuse it.
  Head \val{qthreehead}: $\le$ \val{qthreeheadrate} TFLOP/s, a bigger batch helps ($\ArithmeticIntensity \approx \BatchSize$).
  Conv \val{qthreeconv}: compute bound at \val{qthreeconvrate} TFLOP/s; tune tiles and shapes.}
\ans{4}{Profile the timeline. HBM traffic near \val{hbw} TB/s with kernels back to back: bandwidth, so cut
  bytes. Gaps between short kernels and low traffic: launch overhead, so compile, fuse or capture the graph.}
\end{answers}

\begin{leftcol}
\sheetsec{Worked example}
\begin{practicebox}\boxtitle{Is batch-one decoding limited by math or memory?}
GPT-2 XL (\val{gparams} B parameters, \val{gweightgb} GB at FP16) generates one token at a time on an H100
(\val{hpeaksixteen} TFLOP/s FP16, \val{hbw} TB/s).
\begin{steps}
\item Work and bytes: 2 per parameter $=$ \val{gops} GFLOP on \val{gweightgb} GB, so
  $\ArithmeticIntensity =$ \res{\val{gintensity} FLOP/byte}.
\item Ridge: \val{hpeaksixteen} $\div$ \val{hbw} $=$ \res{\val{ridgesixteen} FLOP/byte}, far above: bandwidth bound.
\item Ceiling: \val{hbw} TB/s $\times$ \val{gintensity} $=$ \val{gattain} TFLOP/s, \res{\val{gshare}\% of peak}.
  Floor: \val{gweightgb} GB $\div$ \val{hbw} TB/s $=$ \res{\val{gfloorms} ms per token}: at most about \val{gtokens} tokens/s.
\end{steps}
\takeaway{Divide, compare with the ridge, then let bytes $\div$ bandwidth set the time.}
\end{practicebox}

\vspace{12pt}
\sheetsec{Your turn: finish it}
\begin{practicebox}\boxtitle{Do INT8 weights help decoding?}
Same model and H100, weights in INT8 (1 byte each), run on the INT8 path (\val{hpeakeight} TOPS).\par\vspace{4pt}
\textbf{Given:} \val{gops} G operations per token on \val{yweightgb} GB; INT8 ridge
\val{hpeakeight} $\div$ \val{hbw} $=$ \val{ridgeeight}.\par
\textbf{You:} find the intensity, the bound, the time floor per token, and say where the speedup comes from.
\end{practicebox}
\end{leftcol}%
\begin{rightcol}
\sheetsec{Common mistakes}
\mistake{``INT8 doubles the peak, so every layer runs twice as fast.''}%
  {A ResNet-50 stage-3 convolution sits at \val{convsixteen} op/byte in FP16. Its activations move
   \val{convact} MB, its weights \val{convweight} MB, so INT8 weights alone lift it to \val{convmixed}, below the
   new ridge \val{ridgeeight}. Only INT8 activations (\val{convall}) clear it.}{}
\mistake{``For batch-one serving, buy the bigger peak.''}%
  {At batch one the ceiling is bandwidth. An A100 has \val{peakratio}$\times$ a T4's peak (\val{apeak} vs.\
   \val{tpeak} TFLOP/s) but \val{bwratio}$\times$ its bandwidth (\val{abw} vs.\ \val{tbw} GB/s) at
   \val{costratio}$\times$ the price: about the same bandwidth-bound throughput per dollar.}{}
\mistake{``Pruning half the weights to zero halves the time.''}%
  {Scattered zeros need index lookups and scattered reads, and can run slower than dense. Only the 2:4 pattern
   on a sparse path skips the work.}{}

\vspace{2pt}
\sheetsec{Check yourself}
\begin{qs}
\item \assumed{} \val{qonefrac}\% of a training step's baseline time runs \val{qonegain}$\times$ faster on a new
  accelerator; \val{qoneserial}\% is data loading. Speedup? Ceiling? What would you fix first?
\item How many 8 KB SRAM reads cost as much energy as one DRAM access? \assumed{} a MAC that streams 3 operands in and
  1 out of DRAM, against a \val{qtwowidth}-wide systolic array that fetches 2 operands per \val{qtwowidth} MACs.
  Charge \val{qtwomac} pJ of arithmetic per MAC. Energy per MAC for each?
\item On the H100 at FP16, place LayerNorm (\val{qthreeln} op/byte), ResNet-50's classifier head at batch
  \val{qthreebatch} (\val{qthreehead}) and the stage-3 convolution (\val{qthreeconv}). Ceiling and lever for each?
\item A batch-one kernel runs far below peak. Is it bandwidth bound or launch bound? What observation decides it,
  and which fix follows each?
\end{qs}
\end{rightcol}
\end{document}
